\chapter{Hybridization thermodynamics}

Two complementary DNA strands are mixed at 100 nM each. Their dissociation
constant under the stated conditions is 5 nM. Predict the duplex concentration.
Is it nearly 100 nM because the strands match perfectly? Is it 95.2 nM
because $100/(100+5)=0.952$? Neither answer respects both strand inventories.
The equilibrium duplex concentration is 80 nM.

We will derive that result, implement a numerically stable solver, and then
connect equilibrium to a restricted sequence-dependent energy model.
Chapter 5 established the orientation of complementary strings. This chapter
adds concentrations, standard states, and temperature. The next chapter will
ask how fast a reaction approaches its predicted equilibrium.

The prerequisites are algebra, logarithms, elementary differentiation, and
Python. No laboratory outcome is claimed by the examples. Our molecular
model contains two distinct single-strand species and one duplex species,
not every structure that real DNA can adopt.

\section{Draw the molecular inventory before the equation}

Call the distinct strands A and B and the duplex D. In this notation,
A names a strand species, not the adenine base. The allowed association is
\begin{equation}
\mathrm{A}+\mathrm{B}\rightleftharpoons\mathrm{D}.
\label{eq:association}
\end{equation}
The reverse arrow matters: an equilibrium mixture can contain both duplex
and unbound strands even when the sequences are perfectly complementary.
It does not mean that the populations stop exchanging molecules.

\Plate{01-inventory}{The three-species model. Each duplex contains one A strand and one B strand, with antiparallel backbones. Solid paths denote strand connectivity; dashed rungs denote pairing. Association arrows change molecular populations, not covalent sequence identity. Hairpins, mismatches, aggregates, and alternative duplexes are outside this inventory.}{fig:inventory}

Let $a$ and $b$ be the total concentrations of A and B strands, measured
in mol/L, and let $x=[D]$ be the duplex concentration. Conservation gives
\begin{equation}
[A]=a-x,\qquad [B]=b-x,\qquad 0\leq x\leq\min(a,b).
\label{eq:conservation}
\end{equation}
A duplex contributes one strand to each total. Consequently, $[D]$ is not
the total concentration of strands residing in duplexes: that quantity is
$2x$. A change of denominator can otherwise make two correct measurements
appear inconsistent.

For the equal-total example, $x=80$ nM leaves 20 nM of each free strand.
The fraction of A bound is $x/a=0.8$, the fraction of B bound is $x/b=0.8$,
and the fraction of all strands in duplexes is $2x/(a+b)=0.8$.
For unequal totals these three fractions need not agree.

\section{Connect free energy to a dimensionless equilibrium constant}

In an ideal dilute solution, approximate each species activity by its
concentration divided by a reference concentration $c^\circ=1$ M.
The dimensionless association constant is
\begin{equation}
K^\circ
=\frac{[D]/c^\circ}{([A]/c^\circ)([B]/c^\circ)}
=\frac{c^\circ[D]}{[A][B]}.
\end{equation}
The concentration-form \Term{dissociation constant} is
$K_d=[A][B]/[D]$, with units of mol/L. Thus
\begin{equation}
K^\circ=\frac{c^\circ}{K_d},\qquad
\Delta G^\circ=-RT\ln K^\circ
=RT\ln\left(\frac{K_d}{c^\circ}\right).
\label{eq:standard}
\end{equation}
Here $R$ is the gas constant and $T$ is absolute temperature.
The logarithm acts on a dimensionless ratio, never on a dimensional number
whose meaning changes when the same concentration is written in nM.

At $T=310.15$ K, take $R=1.98720425864083$ cal/(mol K).
For $K_d=5\times10^{-9}$ M, Eq.~\ref{eq:standard} gives approximately
$-11.78$ kcal/mol. A smaller dissociation constant corresponds to a more
negative standard association free energy.

\subsection{Standard free energy is not the free energy of every mixture}

At arbitrary composition, the reaction free energy is
\begin{equation}
\Delta G=\Delta G^\circ+RT\ln Q,\qquad
Q=\frac{c^\circ[D]}{[A][B]}.
\end{equation}
At equilibrium, $Q=K^\circ$ and $\Delta G=0$. A negative
$\Delta G^\circ$ therefore does not mean the final equilibrium reaction
free energy remains negative. Mixing and depletion contribute through $Q$.

This also explains why complementarity and binding fraction are different
questions. The molecular interaction model determines an equilibrium
constant under chosen conditions; conservation and starting inventories
then determine the populations. Neither layer can substitute for the other.

\section{Solve the finite-pool problem}

Substitute Eq.~\ref{eq:conservation} into the concentration-form constant:
\begin{equation}
K_dx=(a-x)(b-x)
       =ab-(a+b)x+x^2.
\end{equation}
Hence $x$ is a root of
\begin{equation}
x^2-(a+b+K_d)x+ab=0.
\end{equation}
The larger root is outside the physical interval except at a degenerate
complete-binding limit. With $s=a+b+K_d$, the physical root is
\begin{equation}
x=\frac{s-\sqrt{s^2-4ab}}{2}
 =\frac{2ab}{s+\sqrt{s^2-4ab}}.
\label{eq:root}
\end{equation}
The second form follows by multiplying numerator and denominator by
$s+\sqrt{s^2-4ab}$. It avoids subtracting nearly equal large numbers when
binding is weak. We additionally compute
\begin{equation}
s^2-4ab=(a-b)^2+K_d\{2(a+b)+K_d\},
\end{equation}
which is a sum of nonnegative terms.

\Listing{duplex_molar}

The implementation divides all three input concentrations by their maximum
before evaluating the root, then restores concentration units. Scaling limits
unnecessary overflow and underflow in intermediate products. It does not
turn arbitrary floating-point inputs into exact arithmetic: the tests cover
a declared range of concentrations and physical invariants.

Check the opening example in nM units: $s=205$ nM and
$\sqrt{s^2-4ab}=\sqrt{42025-40000}=45$ nM. Thus
$x=(205-45)/2=80$ nM. Substituting back gives
$[A][B]/[D]=20\cdot20/80=5$ nM.
Our public function expects molar inputs, so the corresponding call is
\texttt{duplex\_molar(100e-9, 100e-9, 5e-9)}.

\subsection{An approximation with a missing assumption}

If B is in sufficient excess that its free concentration barely changes,
then $[B]\approx b$. From $K_d=[A][B]/[D]$ and $a=[A]+[D]$ we obtain
\begin{equation}
\frac{x}{a}\approx\frac{b}{K_d+b}.
\label{eq:excess}
\end{equation}
This is a reservoir approximation. At $a=b=100$ nM it predicts 0.952,
but binding would substantially deplete B; the assumption contradicts its
own predicted outcome.

\Plate{03-binding-curve}{Exact finite-pool binding versus the ligand-excess approximation, with A total fixed at 100 nM and $K_d=5$ nM. The approximation fails strongly when B is depleted. The marked equal-total point has 80 nM duplex, not 95.2 nM. Curves are computed from the declared model, not experimental data.}{fig:binding}

For $b=5$ nM, Eq.~\ref{eq:excess} even predicts half of 100 nM A bound,
requiring 50 nM B when only 5 nM exists. Conservation is a powerful
negative control because it catches a plausible-looking curve without
knowing any sequence parameters.

The solver accepts zero inventory and returns zero duplex. It also accepts
$K_d=0$ as the mathematical complete-binding limit, returning $\min(a,b)$;
that is not a claim of an exactly zero measured dissociation constant.
Negative or nonfinite concentrations are rejected.

\section{Separate enthalpy and entropy, then reunite their units}

Within a temperature range where heat-capacity differences are neglected,
\begin{equation}
\Delta G^\circ(T)=\Delta H^\circ-T\Delta S^\circ.
\label{eq:gibbs}
\end{equation}
For duplex formation, the net energetic and configurational contributions
often compete. A favorable enthalpy alone is not the full free-energy change.
Nor should the net entropy change be identified solely with DNA chain
ordering: solvent and ions are part of the physical system too.

\Plate{02-free-energy}{A signed energy balance for the chapter's sequence model at 37 degrees Celsius. A favorable enthalpy and opposing $-T\Delta S^\circ$ term combine into the standard free energy. These are net model terms, not separately measured hydrogen-bond and stacking energies. The actual mixture additionally contributes $RT\ln Q$.}{fig:energy}

We store enthalpy in kcal/mol but entropy and $R$ in cal/(mol K).
The factor 1000 in the implementation is therefore essential:
\begin{equation}
K_d(T)=c^\circ\exp\left[
\frac{1000\Delta H^\circ_{\mathrm{kcal}}-T\Delta S^\circ_{\mathrm{cal}}}
{RT}\right].
\end{equation}

\Listing{kd_from_thermo}

Temperature enters in kelvin, not Celsius. The function rejects nonpositive
temperatures and exponential arguments outside its numerical budget.
It does not silently return an infinity or an exact zero and present that
as a physical prediction.

The error sensitivity is exponential:
\begin{equation}
\frac{K_{d,\mathrm{perturbed}}}{K_d}
=\exp\left(\frac{\delta G^\circ}{RT}\right).
\end{equation}
At 310.15 K, an assumed $+1$ kcal/mol shift multiplies $K_d$ by about 5.07.
This is a sensitivity calculation, not a universal mismatch penalty.

\section{Build a restricted nearest-neighbor model}

An isolated base-pair count ignores sequence order. A
\Term{nearest-neighbor model} assigns contributions to adjacent paired
layers, with additional end terms. SantaLucia's 1998 Table 2 gives a
specific parameter convention at 1 M NaCl \cite{santalucia}.
We use that table consistently rather than mixing initiation terms from
one parameterization with step energies from another.

\Plate{04-neighbors}{A sequence is decomposed into overlapping adjacent duplex steps. AC/TG means upper $5'$--AC--$3'$ above lower $3'$--TG--$5'$; the slash does not put both strands in stored $5'$-to-$3'$ order. End corrections are counted once at each end. Overlapping steps are energy-model terms, not separate molecules.}{fig:neighbors}

The ten representative upper-strand steps and their parameters are:
\begin{center}\small\input{\Generated/parameters.tex}\end{center}
For a step absent from this list, use its reverse complement: AC is the
same duplex-step orientation class as GT. A length-$L$ perfectly matched
duplex has $L-1$ such terms. The terminal contribution at each end is
$(2.3,4.1)$ for A/T and $(0.1,-2.8)$ for G/C, in the table's enthalpy
and entropy units. We reject self-complementary sequences because their
symmetry and species accounting need a separate treatment.

\Listing{nn_parameters}

Walk through our original sequence \texttt{ACGTCAGT}. Its seven upper
steps are AC, CG, GT, TC, CA, AG, GT. Mapping to table representatives gives
three GT terms and one each of CG, GA, CA, and CT. Both ends are A/T.
The enthalpy sum is therefore
\begin{align}
\Delta H^\circ
&=3(-8.4)-10.6-8.2-8.5-7.8+2(2.3)\\
&=-55.7\ \mathrm{kcal/mol}.
\end{align}
The entropy calculation gives $-152.1$ cal/(mol K).
At 310.15 K, $\Delta G^\circ=-8.526185$ kcal/mol and
$K_d\approx981.909$ nM. These values describe our restricted calculation,
not a new experimental measurement.

Compare it with \texttt{AGCTACGT}. Both sequences contain two occurrences
of each base and have the same terminal identities; their neighbor counts differ.
\begin{center}\small\input{\Generated/results-table.tex}\end{center}
Here enthalpy is kcal/mol, entropy is cal/(mol K), $K_d$ is evaluated at
37 degrees Celsius, and $T_m$ uses 100 nM of each strand. Equal GC content
does not force equal model free energies. The approximately 0.65-degree
difference in predicted melting points is not evidence that experiments
could reliably resolve that difference under unspecified conditions.

\section{Derive the concentration term in melting temperature}

For equal totals $a=b=C$, define the \Term{melting temperature} as the
temperature where half of each strand population is duplex-bound:
$x=C/2$. Substitution into mass balance gives
\begin{equation}
K_d(T_m)=\frac{(C/2)^2}{C/2}=\frac C2.
\end{equation}
Combining this with Eq.~\ref{eq:standard} and Eq.~\ref{eq:gibbs},
\begin{equation}
1000\Delta H^\circ
=T_m\left[\Delta S^\circ+R\ln\left(\frac{C}{2c^\circ}\right)\right],
\end{equation}
so
\begin{equation}
T_m=
\frac{1000\Delta H^\circ}
{\Delta S^\circ+R\ln(C/(2c^\circ))}.
\label{eq:tm}
\end{equation}
If $C_T=a+b=2C$ denotes the total of both strand species, the logarithm
contains $C_T/(4c^\circ)$. That familiar factor of four follows from the
concentration definition; it is not an empirical adjustment.

\Listing{melting_kelvin}

\Plate{05-melting}{Calculated duplex fractions for ACGTCAGT and its distinct reverse-complement partner at three equal strand totals. The half-bound line identifies $T_m$; increasing concentration shifts it upward in this model. Constant enthalpy/entropy extrapolation and a three-species inventory are assumptions, not experimentally established validity over the full plotted range.}{fig:melting}

The computed melting temperatures are 19.88, 27.11, and 34.70 degrees Celsius
for 10, 100, and 1000 nM of \emph{each} strand. Tests independently evaluate
$K_d(T_m)$ and the mass-balance root and require a bound fraction of 0.5.
A function returning a plausible temperature is not enough.

For unequal totals, the two strand-specific bound fractions differ.
If $a\geq b$ and melting is defined by half of the limiting B being bound,
$x=b/2$ gives $K_d=a-b/2$. Half of \emph{all} strands bound can even be
unattainable at strongly unequal totals. Always state the fraction being measured.

\section{Know where this model ends}

Salt conditions, mismatch identities, alternative structures, and parameter
versions are part of a physical-model specification. NUPACK's documentation
exposes material, temperature, ensemble, and ionic settings separately
\cite{nupack}. Our implementation intentionally accepts only the declared
1 M NaCl convention. Entering a different sodium concentration raises an
error rather than applying an unreviewed correction.

A mismatch is not modeled by deleting one canonical pair contribution or
adding an arbitrary fixed penalty. It changes the relevant local contexts;
end mismatches, internal mismatches, bulges, and dangling ends need the
appropriate parameter and structural treatment. Similarly, a single-strand
hairpin would consume strand inventory that our three-species model assumes
is available for duplex association.

The two-state approximation groups many molecular configurations into
single-strand and duplex states. A real melting curve can include partially
paired states and competing complexes. Measured absorbance also requires
a signal model and baseline treatment before it becomes a bound fraction.
Our generated curves are population predictions, not simulated instrument traces.

Finally, $K_d$ does not determine a reaction time. In a compatible simple
mass-action kinetic model, $K_d=k_{\mathrm{off}}/k_{\mathrm{on}}$.
Multiplying both rates by the same positive factor leaves equilibrium
unchanged while changing relaxation speed. Chapter 7 develops this distinction
instead of using a low free energy as a synonym for a fast computation.

\section{Reproduce, challenge, and extend}

Run the examples and their regression tests from the repository root:
\begin{Verbatim}[fontsize=\small]
python3 drafts/ch06/reference.py
python3 -m pytest tests/test_ch06_reference.py
\end{Verbatim}
The tests include an independent bisection solution, conservation over
log-spaced concentration ranges, strand-exchange symmetry, fixed expected
step counts, reverse-complement invariance, and melting-point consistency.
The PDF's tables and curves are generated from the same reference results.

The conceptual pipeline is now explicit: oriented sequence, specified
energy model, equilibrium constant, conserved totals, and predicted population.
Each arrow introduces an assumption. Extending the model means revisiting
those assumptions and tests, not merely appending a more complicated formula.

\Needspace{10\baselineskip}
\section{Exercises with worked solutions}

\subsection*{1. Check the inventory}
With $a=100$ nM, $b=30$ nM and $x=20$ nM, find the free concentrations.
\textbf{Solution:} A is 80 nM and B is 10 nM; the implied $K_d$ is 40 nM.
The fractions of A and B bound are 0.2 and $2/3$, respectively.

\subsection*{2. Identify a forbidden prediction}
An approximation predicts 60 nM duplex from 100 nM A and 20 nM B.
\textbf{Solution:} it exceeds the B inventory. No favorable energy can
create the missing strands.

\subsection*{3. Recover the stable root}
Rationalize $(s-\sqrt{s^2-4ab})/2$.
\textbf{Solution:} multiply by the conjugate; the numerator becomes
$s^2-(s^2-4ab)=4ab$, giving Eq.~\ref{eq:root}. This avoids cancellation
in the small root.

\subsection*{4. Audit units}
Why not compute $\ln(5)$ for a dissociation constant of 5 nM?
\textbf{Solution:} the ratio to 1 M is $5\times10^{-9}$, not 5.
Changing display units must not change the physical free energy.

\subsection*{5. Explain equilibrium free energy}
If $\Delta G^\circ<0$, why is $\Delta G=0$ at equilibrium?
\textbf{Solution:} the composition term $RT\ln Q$ compensates the standard
term when $Q=K^\circ$. The standard state is not the final mixture.

\subsection*{6. Count overlapping steps}
How many nearest-neighbor terms occur in an eight-base duplex?
\textbf{Solution:} seven, plus the separately specified end terms.
The overlap shares layers; it does not count seven separate duplex molecules.

\subsection*{7. Verify an orientation class}
Why is upper AC looked up as GT?
\textbf{Solution:} reversing and complementing AC gives GT, representing
the same oriented Watson--Crick duplex step after exchanging the reading direction.
Plain reversal CA is not the same transformation.

\subsection*{8. Recover the factor of four}
If each strand is 100 nM, what $C_T$ belongs in the $C_T/4$ convention?
\textbf{Solution:} $C_T=200$ nM. Its quarter is 50 nM, matching
$K_d(T_m)=C/2$.

\subsection*{9. Inspect the complete-binding limit}
What does $K_d\to0$ imply for unequal totals?
\textbf{Solution:} $x\to\min(a,b)$. Excess strands remain free;
complete binding of the limiting species does not mean all strands are paired.

\subsection*{10. Perturb an energy}
Estimate the $K_d$ multiplier for a $+1$ kcal/mol free-energy perturbation
at 310.15 K. \textbf{Solution:} $\exp[1000/(R\cdot310.15)]\approx5.07$.
The number quantifies sensitivity, not the effect of every mismatch.

\subsection*{11. Distinguish equilibrium and rate}
Two systems have identical $K_d$ but both rate constants in the second
are ten times larger. \textbf{Solution:} their equilibrium populations
match under the same totals; their time evolution generally does not.
An equilibrium calculation cannot determine which reaches the endpoint first.

\subsection*{12. Design a validation study}
How would you challenge this restricted model?
\textbf{Worked rubric:} specify sequences, ionic conditions, strand totals,
temperature protocol, instrument baselines, independent replicates, and
competing-structure checks. Compare predictions with measurements and
uncertainty; do not fit parameters and report validation on the same examples.
Separate equilibrium disagreement from failure to equilibrate.
